
This section outlines the overarching research targets, execution timeline, core deliverables, and academic achievements resulting from the collaborative research project between Nagoya University, ABB Drives Oy, and HighQ Power Oy.

\section*{Research Target and Scope}
The primary objective of this research collaboration is the design, optimization, and thermal management of high-frequency transformers with integrated leakage inductance for high-power DC--DC conversion systems. The overarching system specifications and design priorities are defined as follows:
\begin{singlespace}
    \begin{itemize}
    \item \textbf{Application:} DC/DC conversion systems (with and without galvanic isolation).
    \item \textbf{Power Rating:} $100\text{ kW}\text{--}300\text{ kW}$.
    \item \textbf{DC-Link Voltage Range:} $600\text{ V}\text{--}1500\text{ V}$.
    \item \textbf{Switching Frequency:} $15\text{ kHz}\text{--}30\text{ kHz}$.
    \item \textbf{Target Topologies:} Dual Active Bridge (DAB), LLC/CLLC resonant, phase-shift, and multiphase DC--DC converters.
    \item \textbf{Target Component:} High-frequency isolation transformer with integrated leakage inductance.
    \item \textbf{Core Design Priorities:} High efficiency, high insulation breakdown voltage, compact size and cost reduction, with a primary emphasis on \textbf{advanced thermal management}.
\end{itemize}
\end{singlespace}


\section*{Research Plan and Execution Timeline}
The communication on this collaborative research started in 2024 3Q to 2025 1Q. Initial research scope discussions and preliminary technical training were done in 2024 4Q. The core research and development activities were executed across a structured six-quarter timeline from April 2025 to September 2026:

\begin{itemize}
    \item \textbf{Phase 1: Transformer Design, Modeling, and Comparison (2025 2Q -- 2025 3Q)}
    \begin{itemize}
        \item Development of parametric geometry and magnetics design methodologies, alongside the selection of winding and core loss models.
        \item Proposal of the winding configuration an thermal model.
        \item Proposal of parameters for the evaluation circuit in this study.
        \item validation of the design via 1~kW Finite Element Analysis (FEA) and circuit simulations.
        \item Comparative magnetics and thermal performance evaluation between conventional edgewise and proposed flatwise winding configurations.
    \end{itemize}
    
    \item \textbf{Phase 2: 1~kW Prototype Design and Fabrication (2025 2Q -- 2026 1Q)}
    \begin{itemize}
        \item Selection of core materials, finalization of insulation clearances, transformer terminal design, and cooling environment setup for both flatwise and edgewise prototypes.
        \item Final simulation confirmation of the transformer before fabrication.
        \item Manufacturing lead time and procurement of the two comparative 1~kW transformer prototypes.
        \item Design, fabrication, and functional testing of the 1~kW bidirectional CLLC resonant converter PCB.
    \end{itemize}

    \item \textbf{Phase 3: Experimental Validation and Extended Analysis (2025 4Q -- 2026 2Q)}
    \begin{itemize}
        \item Transformer characterization impedance analyzer measurements.
        \item full-load circuit operational testing, Zero-Voltage Switching (ZVS) verification, and initial thermal performance evaluation.
        \item Extended continuous natural-cooling experimentation until thermal steady-state, accompanied by system efficiency measurements.
        \item Electrostatic FEA simulations to evaluate dielectric breakdown capabilities and electric field stress.
    \end{itemize}

    \item \textbf{Phase 4: Scalability Investigation and Documentation (2026 2Q -- 2026 3Q)}
    \begin{itemize}
        \item Comprehensive scalability analysis of the proposed transformer architecture in terms of both spatial geometry (aspect ratio variation) and power rating ($1\text{ kW}$, $10\text{ kW}$, and $100\text{ kW}$).
        \item Parameter adjustment, 3D modeling, and coupled electromagnetic-thermal FEA simulations across scaled power levels.
        \item Synthesis of final deliverables, including the Master's thesis, summative technical report, and peer-reviewed journal manuscript.
    \end{itemize}
\end{itemize}

\textbf{Final Project Deliverables:}
\begin{enumerate}
    \item A validated, reliable simulation-based design framework for the proposed flatwise-wound transformer concept demonstrated at 1~kW.
    \item Comprehensive design guidelines and performance insights regarding transformer scalability across both \textbf{power rating} and \textbf{physical geometry}.
    \item Dissemination of research findings through a formal Master's thesis and peer-reviewed academic publications.
\end{enumerate}

\section*{Research Achievements}
The collaborative research between Nagoya University, ABB Drives Oy, and HighQ Power Oy yielded the following peer-reviewed publications and academic honors:
\begin{singlespace}
\subsection*{Journal Papers (1)}
\begin{singlespace}
\begin{enumerate}
    \item \textbf{K. Roque}, I. Masuda, S. Choi, J. Imaoka, M. Yamamoto, \textbf{M. Kauhanen}, and \textbf{K. Tashiro}, ``Feasible Evaluation of Edgewise and Flatwise Rectangular Windings for Thermal Management in Transformers,'' \textit{IEEJ Journal of Industry Applications} (Under Review).
\end{enumerate}
\end{singlespace}

\subsection*{International Conference Papers (3)}
\begin{singlespace}
\begin{enumerate}
    \item \textbf{K. Roque}, I. Masuda, S. Choi, J. Imaoka, M. Yamamoto, \textbf{M. Kauhanen}, and \textbf{K. Tashiro}, ``Flatwise Windings for Improved Thermal Management in High Power Transformer Applications,'' in \textit{Proc. 41st Applied Power Electronics Conference (APEC)}, San Antonio, TX, USA, Mar. 22--26, 2026.
    \item \textbf{K. Roque}, I. Masuda, S. Choi, J. Imaoka, M. Yamamoto, \textbf{M. Kauhanen}, and \textbf{K. Tashiro}, ``Comparative Analysis of Flatwise and Edgewise Windings for Improved Thermal Management in Transformers,'' in \textit{Proc. 10th International Power Electronics Conference (IPEC--ECCE Asia)}, Nagasaki, Japan, May 31--Jun. 4, 2026.
    \item \textbf{K. Roque}, I. Masuda, S. Choi, J. Imaoka, M. Yamamoto, \textbf{M. Kauhanen}, and \textbf{K. Tashiro}, ``An Investigation on Flatwise Windings for Improved Thermal Management in Transformers,'' in \textit{Proc. International Conference on Materials and Sustainability (ICMaSS)}, Nagoya University, Nagoya, Japan, Dec. 12--14, 2025.
\end{enumerate}
\end{singlespace}
\subsection*{Awards and Recognitions (2)}
\begin{singlespace}
\begin{itemize}
    \item \textbf{Dialogue Preview Session Presenter (Top 8\% Paper):} Recognized among the top 8\% of accepted papers at the \textit{IEEE Applied Power Electronics Conference (APEC)}, 2026.
    \item \textbf{Best Interim Presentation Award:} Awarded for the Master's Midterm Presentation in the G30 Automotive Engineering Program, Nagoya University.
\end{itemize}
\end{singlespace}
\end{singlespace}

\section*{GitHub Repository for this Research}
The GitHub repository for all the assets of this research are in: \href{https://github.com/AngKapitan/ABB-HighQ-Design-Suite}{ABB-HighQ-Design-Suite GitHub Repository}. This contains important files throughout this research, which include but are not limited to: source codes, simulation files, Fusion360 files, progress reports, and more.

